\section{Proof architecture and dependency diagrams}

This chapter records the architecture of the proof as diagrams.  The
diagrams are schematic: they show dependencies and statuses, not
proof terms.

\subsection{The layer diagram}

\begin{verbatim}
Layer 1: string-flow ledger
  word, validity, weight, molecule, PMI

Layer 2: finite base
  explicit 17 states, first big step

Layer 3: unified core
  candidate parameterisation
  segment exclusions d=1, d=2, d=3, d>=4
  CRT candidate and terminal residue
  final valuation inequality (pending in Lean)

Layer 4: cycle bridge
  cycle word extraction
  QB-8 split
  corrected window bound
  failure window contradiction

Layer 5: final assembly
  no cycle -> unbounded
  IsUnboundedOrbit 7 (pending)
\end{verbatim}

\subsection{File import graph}

\begin{verbatim}
FinitePrefix.lean
  imports S6AuditStage1

UnifiedCoreBridge.lean
  imports FinitePrefix, UnifiedCoreAudit

UnifiedCoreAudit.lean
  imports S6Audit, S6AuditStage1, LteMacro

CycleBridge.lean
  imports BlockAutomaton, FinalStatement, S6AuditStage1,
          PhOne, SurvExAudit, Td0Real

DwdbDiv.lean
  imports CycleBridge, UnifiedCoreAudit

FinalStatement.lean
  imports S6Audit, FBeta
\end{verbatim}

\subsection{Conditional interface graph}

\begin{verbatim}
unifiedCoreClosed (open)
  -> decisiveWindowValuationBoundCorrected_of_unified_core
       (compiled conditional)
  -> failureWindowExistenceOfUnifiedCore
       (compiled conditional)
  -> failureWindowExistence (open)
  -> windowBoundToNoCycle_of_failureWindowExistence
       (compiled conditional)
  -> dwdbDivFinalAssembly (compiled conditional)
  -> five_x_plus_one_diverges_at_7 (pending)
\end{verbatim}

\subsection{Topological theorem order}

\begin{enumerate}[leftmargin=*]
\item word-orbit identity;
\item PMI and PMI-B;
\item finite prefix and first-big-step facts;
\item reset-head predecessor and candidate parameterisation;
\item segment equations and exclusions;
\item CRT candidate, terminal residue, size lower bound;
\item final valuation inequality (pending);
\item cycle-word extraction and QB-8 split;
\item reset block and reachability;
\item corrected window bound (conditional);
\item failure-window contradiction;
\item no-cycle-to-unbounded lemma;
\item final theorem (pending).
\end{enumerate}

\subsection{Status flow}

\begin{verbatim}
compiled facts
  -> compiled facts
  -> unified core (one pending Lean statement)
  -> conditional bridge interfaces
  -> final assembly (pending)
\end{verbatim}

\subsection{What the diagrams do not show}

The diagrams show the dependency structure, not the proof size or the
time taken to check a statement.  A node marked compiled may depend on
hundreds of Mathlib lemmas; the diagrams record only the project-level
dependencies that the manual tracks.

\subsection{Status}

\begin{center}
\small
\begin{tabular}{@{}p{0.60\textwidth}p{0.34\textwidth}@{}}
\toprule
Node & Status \\
\midrule
layer 1--3 mathematics & proven \\
layer 3 Lean final inequality & pending \\
layer 4 cycle bridge & compiled (conditional) \\
layer 5 final assembly & pending \\
\bottomrule
\end{tabular}
\end{center}
